\chapter{Introduction}
\label{ch:introduction}

\section{Motivation}
\label{sec:motivation}

\TODO{The gap in one paragraph: a lecture slide deck is readable by a human and
opaque to software. It has no slide boundaries as data, no notion of a concept,
and nothing addressable to attach a question, a review card or a similarity
search to. Every interactive learning feature therefore has to be authored by
hand, which is why most course material stays static.}

\section{Problem Statement}
\label{sec:problem}

\TODO{State the transformation problem precisely, and note that it is not
merely format conversion: converting a PDF to markdown recovers structure but
produces nothing assessable.}

\thesisfigure{f1-transformation}
  {The transformation this thesis addresses. A single opaque artefact becomes a
   set of independently addressable learning objects, each of which unlocks
   capabilities the source PDF cannot support.}
  {fig:transformation}

\Cref{fig:transformation} frames the whole thesis: the input supports none of
the capabilities on the right, and the distinguishing step is that the derived
objects are \emph{materialised} and reused, not produced on demand and
discarded.

\section{Research Questions}
\label{sec:research-questions}

\TODO{The proposal's original questions were framed for a learning-analytics
thesis (data-profiling metrics, event aggregation) while the title describes a
content-pipeline thesis. Reframing below is pending supervisor confirmation ---
resolve this before writing further, because it determines which chapters carry
the weight.}

\begin{description}
  \item[RQ1] How can heterogeneous lecture materials be automatically
    transformed into structured, interactive educational content using large
    language models, and what architecture supports this reliably?
  \item[RQ2] How do students perceive motivation and engagement in a gamified,
    AI-driven learning platform?
  \item[RQ3] Does interactive, AI-generated content improve perceived learning
    efficiency and usability compared with static slides?
\end{description}

\TODO{Be honest about scope here, in the introduction, not buried in
\cref{ch:evaluation}: RQ1 is answerable by design, implementation and technical
evaluation. RQ2 and RQ3 are \emph{perception} questions that require an
empirical study with participants. No such study was conducted. Either scope
them out with the supervisor's agreement, or present a study design as
contributed future work. Do not imply they were answered.}

\section{Scope and Delimitation}
\label{sec:scope}

\TODO{What is in scope: the ingestion pipeline, the architecture supporting it,
the derived learning features, and a systematic audit of the deployed system.
What is out of scope: controlled learning-outcome measurement, multi-tenant
operation at scale, and formal verification of the anti-hallucination
mechanisms.}

\section{Methodological Note on Sources}
\label{sec:method-note}

\TODO{Short but important, and it earns credit. State that all architectural
claims were verified against source code at revision \thesisCommit{} rather
than taken from project documentation, and say why: a systematic comparison
found twelve documented claims that do not hold in the implementation,
including the repository's own headline feature description. Point forward to
\cref{sec:reality-vs-documented}. This is a methodological choice, and stating
it converts a potential embarrassment into evidence of rigour.}

\section{Structure of this Thesis}
\label{sec:structure}

\TODO{One short paragraph per chapter. Write this last, once the chapters
have settled.}
